%%% SECTION 9: A CANCER PATIENT'S JOURNEY THROUGH PHYSICAL AI %%%

[This section should be approximately 10-12 pages. The major theme is that
A Cancer Patient's Journey, Physical AI \cite{patient-journey-paper} provides
definitive evidence that Physical AI can manage a complete oncology trial
patient journey autonomously across ten stages, with documented safety
outcomes, regulatory compliance, and cost savings. This single-patient
simulation, together with the 24-hour multi-patient simulation, constitutes
pivotal evidence that state of the art AI has the ability to understand the
Physical AI application, proficiently work with the applications, and scale
beyond human and prior technology capabilities.]

[PRIMARY SOURCE: national-platform/patient\_journey/ directory containing three
chunked files: 01\_preamble\_introduction\_stages1to4.tex (lines 1-296),
02\_stages5to8.tex (lines 297-654), and
03\_stages9to10\_discussion\_closing.tex (lines 655-876). All three files must
be used together to build this section.]

\subsection{Ten-Stage Pipeline Architecture}
\label{subsec:journey-pipeline}

[Build from national-platform/patient\_journey/01\_preamble\_introduction\_stages1to4.tex,
covering the introduction and the 10-stage pipeline architecture. Provide an
overview of the complete patient journey from prescreening through closeout.
List all ten stages with brief descriptions.]

[MAJOR THEME: The ten-stage pipeline demonstrates end-to-end Physical AI
capability, not just isolated robotic procedures. This comprehensive scope
differentiates Physical AI oncology trials from prior robotic surgery studies
that focused only on surgical outcomes.]

\subsection{Pre-Trial Stages: Prescreening Through Enrollment}
\label{subsec:journey-pretrial}

[Build from national-platform/patient\_journey/01\_preamble\_introduction\_stages1to4.tex.
Cover Stages 1-4 including patient selection, eligibility screening, informed
consent with Physical AI-specific elements, and enrollment. Reference the
Stage 1 timeline figure and the Patient Selection criteria. Detail the
Regulatory Framework Mapping table showing how each stage maps to the adapted
standards (Adaption: ICH Harmonised Guideline \cite{ich-adapt}, Physical AI
Adaption: 21 CFR Part 50 \cite{cfr50-adapt}, Physical AI Adaption: 21 CFR
Part 312 \cite{cfr312-adapt}).]

[Reference the Three-Perspective Diagram System that provides ASCII
visualizations of each stage from regulatory, clinical, and patient
perspectives.]

\subsection{Treatment Stages: Surgical and Therapeutic Interventions}
\label{subsec:journey-treatment}

[Build from national-platform/patient\_journey/02\_stages5to8.tex. Cover
Stages 5-8 including surgical outcomes, therapeutic interventions, monitoring,
and follow-up. Detail the Stage 5 timeline figure and surgical outcomes data.
Include the Adverse Event Profile documenting safety performance. Reference
Physical AI Unification Standard Level \cite{usl-standard} for the USL
scoring applied during treatment stages.]

[MAJOR THEME: Treatment stages demonstrate that Physical AI systems can
deliver clinical interventions with documented safety outcomes comparable to
or better than human-only delivery. More patients can be treated with more
robots and fewer workers, without compromising quality or safety.]

\subsection{Post-Treatment Stages: Closeout and Analysis}
\label{subsec:journey-closeout}

[Build from national-platform/patient\_journey/02\_stages5to8.tex and
03\_stages9to10\_discussion\_closing.tex. Cover Stages 9-10 including final
closeout metrics, regulatory compliance coverage, and Physical AI ecosystem
data flow. Include the Guidance Documents reference and Final Closeout
Metrics.]

\subsection{Regulatory Compliance Coverage}
\label{subsec:journey-compliance}

[Build from national-platform/patient\_journey/02\_stages5to8.tex on
Regulatory Compliance Coverage. Detail how the patient journey demonstrates
compliance with all three adapted regulatory standards at every stage. Show
the mapping between journey stages and specific regulatory provisions.]

\subsection{Physical AI Ecosystem and Data Flow}
\label{subsec:journey-ecosystem}

[Build from national-platform/patient\_journey/02\_stages5to8.tex on the
Physical AI Ecosystem and Data Flow sections. Detail how data flows between
Physical AI systems, clinical staff, regulatory systems, and patients
throughout the ten-stage journey. Reference Physical AI National MCP Servers
\cite{national-mcp-paper} and the MCP protocol \cite{mcp-protocol} as the
infrastructure supporting this data flow.]

\subsection{FDA Cost-Savings Analysis}
\label{subsec:journey-cost}

[Build from national-platform/patient\_journey/02\_stages5to8.tex, FDA
Cost-Savings Results section. Present the detailed financial analysis from
the patient journey, including per-patient cost reductions, time savings at
each stage, and projected savings at trial scale. Reference Tufts CSDD data
\cite{tufts-delay} for context on the value of reducing drug development
timelines. Reference NCI information on trial costs \cite{nci-trials-paying}.]

[MAJOR THEME: The financial benefits of Physical AI are substantial and
documented. Physical AI costs less per patient while providing faster treatment,
more thorough monitoring, and better documentation. These cost savings compound
when scaled to the 168-patient level demonstrated in the 24-hour simulation
and further to nationwide deployment.]

\subsection{Trial Timeline Comparison}
\label{subsec:journey-timeline}

[Build from national-platform/patient\_journey/02\_stages5to8.tex, Trial
Timeline section. Compare the Physical AI trial timeline to traditional trial
timelines. Show where Physical AI enables compression of trial phases and
where it introduces new phases (such as Phase 0 simulation validation). The
net effect should demonstrate significant overall time savings.]

\subsection{Discussion of Patient Journey Evidence}
\label{subsec:journey-discussion}

[Build from national-platform/patient\_journey/03\_stages9to10\_discussion\_closing.tex,
Section 4 (Discussion). Cover autonomous workflow significance, regulatory
adaptations connection, clinical implications, diagram evolution, safety
architecture, and comparison with traditional trials. Synthesize how this
evidence supports the broader National Platform goals.]

[MAJOR THEME: Both the single-patient journey and the 24-hour simulation
are pivotal completed examples towards state of the art AI ability to
understand, work with, and scale Physical AI applications. The evidence
is comprehensive: safety data, regulatory compliance, cost analysis, and
patient outcomes all support the transition to Physical AI oncology trials.]

\subsection{Limitations and Future Directions}
\label{subsec:journey-limitations}

[Build from national-platform/patient\_journey/03\_stages9to10\_discussion\_closing.tex,
Section 5 (Limitations and Future Work). Cover the acknowledged limitations
of the single-patient simulation and how they are addressed by the multi-patient
24-hour simulation and by this National Platform's comprehensive framework.]
